\chapter{Derivadas funcionales típicas en QFT \tiny{(quantum field theory)}}

	\begin{tikzpicture}
	\fill [left color=red!50, right color=teal!50] (0,0) rectangle (6.5,.1);
	\fill [left color=teal!50, right color=blue!50] (6.5,0) rectangle (11.5,.1);
	\end{tikzpicture}

\vspace{10mm}
\begin{adjustwidth}{50pt}{50pt}
\begin{ejemplo}
\vspace{2mm}
Este capítulo va a ser un mini tutorial acerca de la derivada funcional con funciones generalizadas.
\vspace{2mm}

\end{ejemplo}
\end{adjustwidth}

\vspace{5mm}

\section{Derivadas funcionales típicas en QFT}

\begin{definition}
	
	Funcional: $\ F[\phi(x)] = $ número, real o complejo. Derivada funcional (el resultado es una función generalizada, con $u(x)$ una función test):
	
	\vspace{3mm}
	\begin{large}
	\begin{equation}
	\label{T37DF}
	\boldsymbol{
	\fdv{F[\phi(x)]}{\phi(x)} \ = \ \lim_{\varepsilon \to 0} \left\{  \ F \ [ \ \phi(x)\ + \ \varepsilon \ u(x)   \ ] \ - \ F[ \ \phi(x) \ ] \ \right\}
	}	
	\end{equation}
	\end{large}
\end{definition}

\vspace{5mm}
\section{Ejemplos. Propiedades}





\begin{large}{\underline{\textbf{Ejemplo 1}}}\end{large} $\triangleright \qquad F[\phi(x)]=\displaystyle \int_{-\infty}^{+\infty} \phi^2(x)\ \dd x$

$\triangleright \qquad F[\phi(x)]=\displaystyle \int_{-\infty}^{+\infty} \phi^2(x)\ \dd x$

\vspace{2mm} $\displaystyle \fdv{F[\phi(x)]}{\phi(x)}=\ \lim_{\varepsilon \to 0} \ \dfrac 1 \varepsilon \left\{ \ \int_{-\infty}^{+\infty} (\ \phi(x) \ + \ \varepsilon \ u(x) \ )^2 \ \dd x - \int_{-\infty}^{+\infty} \phi^2(x) \ \dd x \ \right\} = $

$\displaystyle = \lim_{\varepsilon \to 0} \ \dfrac 1 \varepsilon \int_{-\infty}^{+\infty} \left( \cancel{\phi^2} + \varepsilon^2 u^2 + \varepsilon 2 \phi u - \cancel{\phi^2} \right) \ \dd x = \lim_{\varepsilon \to 0} \  \int_{-\infty}^{+\infty} (\cancelto{0}{\varepsilon u^2} + 2\phi u)\ \dd x = $

$\displaystyle = 2 \int_{-\infty}^{+\infty}\phi(x)\ u(x9\ \dd x \ \Rightarrow \quad \left(\  \fdv{F[\phi(x)]}{\phi(x)} \ \right)_{\ u}\ = \ \left( \ 2\ \phi(x)\ \right)_{\ u} $

\vspace{2mm} En los textos asuele aparecer sin el subíndice {u} que hace referencia a que se trate de funciones generalizadas y aparece esta propiedad como:

\begin{equation}
\boldsymbol{ \subrayado{ \ 
 F[\phi(x)]=\displaystyle \int_{-\infty}^{+\infty} \phi^2(x)\ \dd x \qquad   \rightarrow \qquad  \fdv{F[\phi(x)]}{\phi(x)} \ = \  2\ \phi(x)
\ } }	
\end{equation}

\vspace{5mm} Nos preguntamos: 
$\displaystyle \quad \text{\Large{?`}} \ \ 
 F[\phi(x)]=\displaystyle \int_{-\infty}^{+\infty} [\phi^2(x)+4\phi^3(x) ] \ \dd x \ \   \rightarrow \ \  \fdv{F[\phi(x)]}{\phi(x)} \ = \  2\ \phi(x)+12\phi^2(x) \ \ \text{\Large{?}} $
 
 La respuesta es \textbf{sí}
 
 \vspace{5mm}
 \begin{myblock}{Propiedad 1 derivadas funcionales}
 \begin{equation}
 \label{T37P1DF}
 \boldsymbol{
 F[\phi(x)] \ = \ \int_{-\infty}^{+\infty} G(\phi(x)) \ \dd x \quad \Rightarrow \quad \fdv{F[\phi(x)]}{\phi(x)} \ = \ \pdv{G(\phi)}{\phi}
 }	
 \end{equation}
	
 \end{myblock}

\vspace{10mm}
\begin{large}{\underline{\textbf{Ejemplo 2}}}\end{large} $\triangleright \qquad F[\phi(x)]=\ \phi(y)$

\vspace{2mm} Por ejemplo, $\ F[x^2]=y^2;\ \ F[\sin x]=\sin y;\ $ etc	.

\vspace{2mm}  $\displaystyle \fdv{F[\phi(x)]}{\phi(x)}=\ \lim_{\varepsilon \to 0} \ \dfrac 1 \varepsilon \left\{ \ \int_{-\infty}^{+\infty} (\cancel{\phi(y)}+\varepsilon u(y) - \cancel{\phi(y)} \ \dd x \ \right\} = u(y) = \int \delta(x-y) \ u(x)\ \dd x$

Luego, $\ \displaystyle F[\phi(x)]=\ \phi(y) \quad \Rightarrow \quad \left( \ \fdv{F[\phi(x)]}{\phi(x)} \ \right)_{\ u}= \left( \ \delta(x-y) \ \right)_{\ u}$

O, como suele aparecer en los libros de texto (sin el subíndice $u$), como $\ F[\phi(x)]=\phi(y)$

\vspace{5mm}
\begin{myblock}{Propiedad 2 derivadas funcionales}
\begin{equation}
 \label{T37P2DF}
 \boldsymbol{
 F[\phi(x)] \ = \  \phi(y) \quad \Rightarrow \quad \fdv{\phi(y)}{\phi(x)} \ = \ \delta(x-y)
 }	
 \end{equation}
\end{myblock}	


\vspace{10mm}
\begin{large}{\underline{\textbf{Ejemplo 3}}}\end{large} $\triangleright \qquad F[\phi(x),\phi(x)]=\ \displaystyle \int \phi\ \psi^2 \ \dd x$	$\qquad $ Ahora se puede calcular tanto $\displaystyle \fdv{F}{\phi}$ como $\displaystyle \fdv{F}{\psi}$


\vspace{2mm}  $\displaystyle \fdv{F}{\phi}=\ \lim_{\varepsilon \to 0} \ \dfrac 1 \varepsilon \left\{ \ \int_{-\infty}^{+\infty} [\ (\phi+\varepsilon u) \psi^2 \ - \ \phi \psi^2 \ ] \right\} = \cdots = \int \psi^2 u(x) \dd x \ \to \ \fdv{F}{\phi}=\psi^2$

Que, como vemos, sigue la regla del ejemplo 1, $\ \ G(\phi,\psi)=\phi\psi^2 \ \ \to\  \ \displaystyle \fdv{G}{\phi}=\psi^2\ $ y, del mismo modo, $\ \ G(\phi,\psi)=\phi\psi^2 \ \ \to\  \ \displaystyle \fdv{G}{\psi}=2\phi \psi$ 





\vspace{10mm}
\begin{large}{\underline{\textbf{Ejemplo 4}}}\end{large} $\triangleright \qquad F[\phi(x)]=\ \displaystyle \int \phi'(x) \ \dd x$

\vspace{2mm} Vamos a ver como tratar con funcionales definidos a partir de derivadas.

\vspace{2mm} 	$\displaystyle \fdv{F}{\phi}=\ \lim_{\varepsilon \to 0} \ \dfrac 1 \varepsilon \left\{ \ \int_{-\infty}^{+\infty}  (\, \phi(x)+\varepsilon u(x) \, )' \ \dd x
\ - \ \ \int_{-\infty}^{+\infty} \phi'(x)\ \dd x
 \right\} =
 \lim_{\varepsilon \to 0} \ \dfrac 1 \varepsilon \int \varepsilon u'(x)\ \dd x=$
 
 $\displaystyle = \int u'(x)\ \dd x = \eval{u(x)}_{-\infty}^{+\infty} \begin{matrix} \\ = \\ \text{\tiny{f. test}} \end{matrix}
 0$

\vspace{2mm}

$$\subrayado{\  \boldsymbol{
 F[\phi(x)] \ = \ \int  \phi'(x)\ \dd x  \quad \Rightarrow \quad \fdv{F[\phi(x)]}{\phi(x)} \ = \ 0
 }	\ }$$
	
 
 \vspace{5mm} Sea ahora, $\ F[\phi(x)]=\displaystyle \int (\phi')^2 \ \dd x$
 
\vspace{2mm} 	$\displaystyle \fdv{F}{\phi}=\ \lim_{\varepsilon \to 0} \ \dfrac 1 \varepsilon \ \int \left\{ \ \left[ \ (\phi+\varepsilon u)' \ \right]^2 \ - \ (\phi')^2   \right\}= 
\lim_{\varepsilon \to 0} \ \dfrac 1 \varepsilon \ \int \left\{ \  \ (\phi'+\varepsilon u')^2  \ - \ (\phi')^2   \right\} =$

$\displaystyle \lim_{\varepsilon \to 0} \ \dfrac 1 \varepsilon  \int (\, \cancel{\phi'^2} +\varepsilon^2 u'^2+2\varepsilon \phi' u' -\cancel{\phi'^2}\,) = \lim_{\varepsilon \to 0} \ \dfrac 1 \varepsilon \int (\, \cancelto{0}{\varepsilon u'^2} \ + \ 2\varepsilon \phi' u' \, ) =$

$\displaystyle = \int 2 \phi'(x) u'(x) \dd x \ 
\begin{matrix} \\ = \\ \text{\tiny{partes test}} \\ \text{\tiny{ec. (\ref{T36partes-test})}} \end{matrix}   \ 
-\int 2\phi'' \ u(x)\ \dd x\, . \quad$ Como función generalizada: $\displaystyle \ \ \fdv{F}{\phi} = -2\phi''$ 

$$\subrayado{\ \boldsymbol{ F[\phi(x)]=\displaystyle \int (\phi')^2 \ \dd x \ \ \to \ \ \fdv{F}{\phi} = -2\phi'' } \ }$$

\vspace{2mm} Es usual que en TCC las acciones, lagrangianos y hamiltonianos dependan de potencias de primeras derivadas.


\vspace{5mm} Generalizando, $\ \ F[\phi(x)]=\ \displaystyle \int (\phi')^n\ \dd x$

\vspace{2mm} \vspace{2mm} 	$\displaystyle \fdv{F}{\phi}=\ \lim_{\varepsilon \to 0} \ \dfrac 1 \varepsilon \ \int [\, (\phi'+\varepsilon u')^n-(\phi')^n \, ] \dd x =  $

$\displaystyle = \lim_{\varepsilon \to 0} \ \dfrac 1 \varepsilon \ \int \left[\, 
\cancel{\mqty(n\\0) (\phi')^n} + \mqty(n\\1) (\phi)^{n-1} \varepsilon u' + \mathcal O(\varepsilon^2) 
-\cancel{\mqty(n\\0) (\phi')^n}
\, \right] \dd x=
\int n (\phi'(x))^{n-1} \ u'(x) \ \dd x \ $

$\displaystyle \
\begin{matrix} \\ = \\ \text{\tiny{partes test}} \\ \text{\tiny{ec. (\ref{T36partes-test})}} \end{matrix}  \ 
\int - n \left( \phi'(x))^{n-1} \right)' \ \dd x = \int - n \, (n-1)\,  (\phi')^{n-2} \, u(x)\ \dd x =$

luego, $\quad \displaystyle F[\phi]=(\phi')^n \ \ \to \ \ \fdv{F}{\phi}\ = \ - n\, (n-1)\, (\phi')^{n-2}\, \phi''\, , \ $ que lo expresamos mejor como:

$$\displaystyle \boldsymbol{ \subrayado{\ 
F[\phi(x)]=\ \displaystyle \int (\phi')^n\ \dd x
\quad \Rightarrow \quad 
\fdv{F}{\phi} \ = \ -\pdv{x} \left[ \pdv{(\phi')^n}{\phi'} \right]
\ }}$$

\vspace{2mm} De modo más general si cabe,

\vspace{5mm}
\begin{myblock}{Propiedad 3 derivadas funcionales}
\begin{equation}
 \label{T37P2DF}
 \boldsymbol{
F[\phi(x)]=\ \displaystyle \int M(\phi') \ \dd x
\quad \Rightarrow \quad 
\fdv{F}{\phi} \ = \ -\pdv{x} \left[ \pdv{M(\phi')}{\phi'} \right]
 }	
 \end{equation}
\end{myblock}	

\vspace{5mm} Comprobemos algunos casos particulares que nos han llevado hasta aquí:

\hspace{1cm}$\centerdot \quad \displaystyle M(\phi')=\phi' \quad \to \quad \pdv{M}{\phi}=1 \quad \to \quad -\pdv{x}(1)=0$ \hspace{4.4cm}$\Box$

\hspace{1cm}$\centerdot \quad \displaystyle M(\phi')=\phi'^2 \quad \to \quad \pdv{M}{\phi}=2\phi' \quad \to \quad -\pdv{x}(2\phi')=-2\phi''$ \hspace{3cm}$\Box$


\vspace{10mm}
\begin{large}{\underline{\textbf{Ejemplo 5}}}\end{large} $\triangleright \qquad F[\phi(x)]=\ \displaystyle \int k (\phi(x),\phi'(x))\ \dd x$	 $\qquad \qquad$ Generalización de las props. anteriores:

\vspace{5mm}
\begin{myblock}{Propiedad 4 derivadas funcionales}
\begin{equation}
 \label{T37P2DF}
 \boldsymbol{
F[\phi(x)]=\ \displaystyle \int K \left(\phi(x),\phi'(x) \right)\ \dd x
\quad \Rightarrow \quad 
\fdv{F}{\phi} \ = \pdv{K}{\phi} \ - \ \pdv{x} \left[ \pdv{K}{\phi'} \right]
 }	
 \end{equation}
\end{myblock}

Ecuación completamente similar a las ecuaciones de Euler-Lagrange.


\vspace{10mm}
\begin{large}{\underline{\textbf{Ejemplo 6}}}\end{large} $\triangleright \qquad F[\psi(x),\phi(x)]=\ \displaystyle \int \psi(x)\ \phi(x)\ \dd x   \ \begin{matrix} \\ = \\ \text{\tiny{x vble. muda}} \end{matrix} \ \int \psi(y)\ \phi(y)\ \dd y  $	

Nos preguntamos. ?`Habrá diferencia entre $\ \displaystyle \fdv{F}{\phi(x)}\ $ y $\ \displaystyle \fdv{F}{\phi(y)}\ $?

\vspace{2mm} ------ $\displaystyle \quad  \fdv{F}{\phi(y)} = \fdv{\phi(y)} \int \psi(y)\phi(y)\ \dd y 
 \ \begin{matrix} \\ = \\ \text{\tiny{ec. \ref{T37P2DF}}} \end{matrix} \ $
 \begin{small}
$\pdv{\phi} \left( \int \psi(y)\phi(y)\ \dd y \right) - \pdv{y} \cancelto{0}{\left( \fdv{\phi'} \int \psi(y)\phi(y)\ \dd y \right) }$ \end{small} 
$= \psi(y)$

\vspace{2mm} ------ $\displaystyle \quad  \fdv{F}{\phi(x)} =  \fdv{\phi(x)}  \int \psi(y)\phi(y)\ \dd y 
\ \begin{matrix} \\ = \\ \text{\tiny{x vble. muda}} \end{matrix} \
  \fdv{\phi(x)}  \int \psi(x)\phi(x)\ \dd x = \psi(x)$
  

Parece raro, pero es así.  La expresión $\ \displaystyle \fdv{F}{\phi(x)} = \fdv{\phi(y)}{\phi(x)}\  \fdv{F}{\phi(y)}\ $ no es correcta, falta una integral, la expresión correcta es: $\ \displaystyle \fdv{F[\phi(y)]}{\phi(x)} = \int \dd y\ \fdv{\phi(y)}{\phi(x)}\  \fdv{F}{\phi(y)}\ $

\color{gris} \underline{Explicación}:

\begin{multicols}{2}
\begin{small}
Para una función de muchas variable, $f(y_1, y_2, \cdots, y_n) \ \to \  \displaystyle \pdv{f}{x_i}=\sum_{j=1}^n \ \pdv{y_j}{x_i} \ \pdv{f}{y_j} \ $ (regla de la cadena para derivadas parciales), un funcional se puede interpretar como que depende de muchas variables ($\phi(x_1), \phi(x_2), \cdots $) y el sumatorio se convierte en integral al pasar al continuo.\end{small}

\begin{figure}[H]
	\centering
	\includegraphics[width=.35\textwidth]{imagenes/img37-01.png}
\end{figure}
\end{multicols}
\color{black}

Tendremos $\displaystyle \qquad \fdv{F[\phi(y)]}{\phi(x)} \ = \ \int \dd y \ \fdv{F[\phi(y)]}{\phi(y)} \ \cancelto{\delta(y-x)}{\fdv{\phi(y)}{\phi(x)}}$

\textcolor{gris}{La delta de Dirac es \emph{par}, $\ \delta(-z)=\delta(z) \ \leftrightarrow \ \delta(y-x)=\delta(x-y)$}

$\displaystyle  \fdv{F[\phi(y)]}{\phi(x)} \ = \ \int \dd y \ \fdv{F[\phi(y)]}{\phi(y)} \ \delta(y-x) $

\vspace{2mm} Generalizando, 

\vspace{5mm}
\begin{myblock}{Propiedad 5 derivadas funcionales}
\begin{equation}
\begin{split}
\boldsymbol{
F[\phi(x)]= \displaystyle \int K \left(\phi(y),\phi'(y) \right)\ \dd y }
& \\boldsymbol{ \Rightarrow } \  \\
& \boldsymbol{\fdv{F[\phi(y)]}{\phi(x)} \ = \int \dd y \ \left[ \pdv{K}{\phi(y)} - \pdv{y} \left( \pdv{K}{\phi'(y)}\right) \right]\ \delta(y-x) } \ \\ 
& \boldsymbol{(\text{\scriptsize{aplicando} } \delta \text{\scriptsize{ Dirac}} ) \ \ =\pdv{K}{\phi(x)} \ - \ \pdv{x} \left[ \pdv{K}{\phi'(x)} \right]	  } 
 \end{split}
 \end{equation}
\end{myblock}




































